% Source: https://artofproblemsolving.com/community/c3h2571226_mathematics
% Stable identity: aops-topic-2571226-q01
Two sides do not determine the third side, since the angle between them can vary. Let the third side be $c$.

\begin{asy}
usepackage("amsmath");
size(5cm);
defaultpen(fontsize(10pt));
pair A=(0,0),B=(13,0),C=7.5*dir(60);
draw(A--B--C--cycle,linewidth(0.9));label("$13$",(A+B)/2,S);label("$7.5$",(A+C)/2,NW);label("$c$",(B+C)/2,NE);
label("$A$",A,SW);label("$B$",B,SE);label("$C$",C,N);
\end{asy}

By the triangle inequality, each side is less than the sum of the other two:
\[c<7.5+13=20.5,\qquad 13<7.5+c\iff c>5.5,\qquad 7.5<13+c\ \text{(always true)}.\]
Conversely, for every $c$ with $5.5<c<20.5$ the three strict inequalities hold, and three lengths satisfying them always form a triangle. So the third side can be any length with
\[\boxed{5.5<c<20.5}.\]
